% ── background.tex · §2 · STATUS: skeleton (expand) ──────────────────
% PP + trace/replay mechanism + λΠ-calculus modulo rewriting. Sets up §4.
\section{Background}\label{sec:background}

\subsection{The Predicate Prover and its proofs}
\pp{} is \ab's automated prover for a typed first-order set theory: given a
proof obligation---a sequent \(H \seq G\) of hypotheses and a goal---it
searches for a proof using a fixed set of inference rules (roughly
\numRulesSpec{} of them, catalogued in \pp's specification: conjunction,
disjunction, implication, equivalence, negation, universal and
existential quantification, equality, normalisation, arithmetic, and
Booleans). Each rule transforms the current sequent, and a successful
search yields a tree of rule applications closing the goal.
% TODO(cite): PP specification / Atelier B reference manual.

\subsection{Traces and replays}\label{sec:replay-format}
\pp{} can be instrumented to record a compact \emph{trace}: a postorder
listing of the rules it applied. \ab's REPLAY utility re-runs a trace and
expands it into a \emph{replay}, in which the main sequent proof is written
\emph{prefix}-style (a rule before its child subproofs) and every line
carries the per-rule formula annotation \pp{} saw when it applied that
rule. Three details matter for reconstruction
(\autoref{sec:replay}): some lines are \emph{phantom} normalisation
markers that carry no inference and must be filtered; the result-chain
child of a branching quantifier is emitted \emph{before} its branch rule;
and the rule name encodes its arity and variant
(\autoref{sec:encoding}). The first non-phantom annotation is the overall
goal.

\subsection{The \texorpdfstring{$\lambda\Pi$}{lambda-Pi}-calculus modulo rewriting}
\lp~\cite{hondet:hal-02981561} and \dedukti~\cite{assaf2023dedukti}
implement the \lpcmr~\cite{CousineauDowek}: a dependently typed
\(\lambda\)-calculus in which the conversion relation used by type-checking
is extended with user-declared rewrite rules. This combination gives a
small, well-understood proof-checking kernel together with the ability to
embed an object logic by declaring its connectives and computation rules,
which is why it is used as an interoperability target for many proof
systems. We exploit both features: \bbook{} connectives and arithmetic are
declared in \texttt{B.lp}, and \pp's structural conventions (\(n\)-ary
conjunction, tuple quantifiers) are made to compute via rewrite rules so
that one encoded rule matches goals of any arity.
